%=============================================================================
% DISCUSSION AND CONCLUSION -- about a third of a page.
%=============================================================================
\section{Discussion}
\label{sec:discussion}

\textbf{Why \model{} works.} Weekly dengue counts are close to a random walk, so the
most informative feature is the last value. A direct head must carry that value through
the encoder, and three published architectures, attention-based, recurrent and
diffusion-based, lose it. \model{} hands the level to the output, so the encoder only
learns the departure, and the gate stops a weak departure from pulling the forecast
away from persistence. This is why one parameter repairs three different designs.

\textbf{What it cannot do.} The anchor removes failure but cannot create information.
The working encoders, the SEIR-LSTM and $k$-NN, which has no network at all, leave
residuals that correlate 0.91--0.97 with one another: the models make the same errors,
and those errors are lags on large moves. When the next move is not in the inputs, as
in the 2026 outbreak, every model lags it, and on later weeks no learned model beats
persistence on RMSE.

\textbf{Physics.} The SEIR simulator cannot be the forecaster, because 14--16\% of
targets lie below its zero-transmission floor and its force of infection is poorly
predictable. Inside the gated head it looked useful, but the ablation and the
prospective test both show that the anchor and gate do the work. For this target, the
useful prior is persistence, not transmission dynamics. We therefore recommend that
dengue forecasting studies report persistence, a matched no-physics control and, where
possible, a prospective test.

\textbf{Limitations.} One country, one disease and district-week resolution; three
origins measure run-to-run stability more than the series; the prospective test is
amended, covers one 2.4-year period dominated by one outbreak, and uses one encoder
family; and the graph is the benchmark's hand-built adjacency rather than mobility.

\section{Conclusion}
We proposed \model{}, a persistence-anchored gated head that plugs into any
spatio-temporal encoder at the cost of one parameter, together with a leakage-free
Sri Lankan dengue dataset rebuilt from source reports and extended to 2026. \model{}
turns three published ST-GNNs that are worse than persistence into usable forecasters,
and with a strong encoder it beats persistence on every paired development run. The
ablation shows that the gain comes from the anchor and the gate, not from the SEIR
simulator. On 2024--2026 weeks no learned model beats persistence on RMSE, although
\model{} is better on MAE and SMAPE. Future gains must come from inputs that anticipate
outbreaks, not from the architecture.
